\documentclass[a4paper,10pt]{article}

\usepackage[utf8]{inputenc}
\usepackage[T1]{fontenc}
\usepackage{amsmath}
\usepackage{amssymb}
\usepackage{mathtools}
\usepackage[pdfpagemode=UseNone]{hyperref}

\DeclareMathOperator{\SD}{SD}
\DeclareMathOperator{\Var}{Var}
\DeclarePairedDelimiter\bracket{[}{]}
\newcommand{\E}[1]{\mathbb{E}\bracket*{#1}}

\begin{document}

% --------------------------------------------------------------------------------
\section*{Definitions}

As before, we have the following quantities:

\begin{verbatim}
S = The set that contains all legal chess positions as a subset
N = The size of S (2.19645e53)
S_sample = A random sample from S
N1 = The size of S_sample (2.746e15)
S_n1_candidates = The computed subset of S_sample containing
                  all potentially legal positions
N_n1_candidates = The size of S_n1_candidates (6030640)
S_2 = A random sample from S_n1_candidates
N2 = The size of S_2 (1e5)
N_s2_legal = The number of legal positions in S_2 (99983)
N_s2_illegal = The number of illegal positions in S_2 (17)
S_n1_legal = The set of legal positions in S_sample
N_n1_legal = The size of S_n1_legal (not computed exactly)
N_n1_legal_est = Estimated number of legal positions
                 in S_sample (6029615)
N_legal_est = Estimated number of legal chess positions
              (4.823692e44)
\end{verbatim}

We introduce the following unknown parameters:

\begin{itemize}

  \item $N_{\text{legal}}$ = The number of positions in $S$ that are legal chess
    positions. $S_{\text{legal}}$ is the corresponding set. This is the main
    unknown parameter to be estimated, $\approx 4.8 \cdot 10^{44}$.

  \item $N_{\text{candidates}}$ = The number of positions in $S$ that are
    potentially legal according to the first automatic analysis.
    $S_{\text{candidates}}$ is the corresponding set. This is the nuisance
    parameter, $\approx 4.8 \cdot 10^{44}, > N_{\text{legal}}$.

\end{itemize}

Introduce the following variables:

\begin{align}
n &= N_1                                                          \\
m &= N_2                                                          \\
C &= N_{\text{n1\_candidates}}                                    \\
L &= N_{\text{n1\_legal}}                                         \\
H &= N_{\text{s2\_legal}}                                         \\
p &= \frac{N_{\text{candidates}}}{N}                              \\
q &= \frac{N_{\text{legal}}}{N_{\text{candidates}}}
\end{align}

That is, $n$ is the size of the first sample, $m$ is the size of the second
sample, $p$ is the fraction of positions in $S$ that are potentially
legal and $q$ is the fraction of positions in $S_{\text{candidates}}$ that are legal.

Since $p$ can be estimated by $\frac{C}{n}$ and $q$ by $\frac{H}{m}$, the
estimate $\hat{N}_{\text{legal}}$ of the number of legal chess positions is:

\begin{equation}
  \hat{N}_{\text{legal}} = N \frac{C}{n} \frac{H}{m}
\end{equation}

% --------------------------------------------------------------------------------
\section*{Expected value}

Because of the way the sampling procedure is defined, $C$, $L$ and $H$ are
random variables with the following distributions:

\begin{align}
  C & \sim \operatorname{Binomial}(n, p)                     \\
  L \mid C & \sim \operatorname{Binomial}(C, q)              \\
  H \mid C,L & \sim \operatorname{HyperGeometric}(C, L, m)
\end{align}

For the distributions, we have~\cite{wiki:binomial,wiki:hypergeom}:

\begin{align}
  \E{C} &= np                           \\
  \E{L \mid C} &= Cq                    \\
  \E{H \mid C,L} &= m \frac{L}{C}
\end{align}

so

\begin{align}
  \E{\hat{N}_{\text{legal}} \mid C,L} & = \E{N \frac{C}{n} \frac{H}{m} \mid C,L}   \\
    &= \frac{N}{n} \frac{C}{m} \E{H \mid C,L}                                      \\
    &= \frac{N}{n} \frac{C}{m} m \frac{L}{C} = \frac{N}{n} L
\end{align}

Taking expectations again and using the law of iterated
expectations~\cite{wiki:totalexpect},

\begin{equation}
  \E{\hat{N}_{\text{legal}} \mid C} = \frac{N}{n} \E{L \mid C} = \frac{N}{n} Cq
\end{equation}

Taking expectations once more,

\begin{align}
  \E{\hat{N}_{\text{legal}}} &= \frac{N}{n} q \E{C} = \frac{N}{n} qnp = Npq              \\
    &= N \frac{N_{\text{candidates}}}{N} \frac{N_{\text{legal}}}{N_{\text{candidates}}}  \\
    &= N_{\text{legal}} \label{eq:unbiased}
\end{align}

Therefore the estimator is unbiased.

% --------------------------------------------------------------------------------
\section*{Variance}

The variance of a product of two independent random variables $X$, $Y$,
is~\cite{wiki:productvariance}:

\begin{equation}
  \Var(XY) = \E{X}^2\Var(Y) + \E{Y}^2\Var(X) + \Var(X)\Var(Y)
\end{equation}

which can be written

\begin{equation}
  \label{eq:rel.var}
  \frac{\Var(XY)}{\E{XY}^2} = a + b + ab
\end{equation}

where

\begin{equation}
  \label{eq:rel.ab}
  a = \frac{\Var(X)}{\E{X}^2}, \qquad b = \frac{\Var(Y)}{\E{Y}^2}
\end{equation}

The estimator is

\begin{equation}
  \label{eq:nhat}
  \hat{N}_{\text{legal}} = N \frac{C}{n} \frac{H}{m}
\end{equation}

Let

\begin{equation}
  X = \frac{C}{n}, \qquad Y = \frac{H}{m} \label{eq:xy}
\end{equation}

Although H is obtained by sampling from the C candidates, the legality of each
candidate is an independent Bernoulli~\cite{wiki:bernoulli} random variable with
probability q, so\footnote{
  Assuming $C \ge m$. For sample sizes of interest, $\Pr(C < m)$ is extremely
  small. If $C < m$, all $C$ candidates can instead be analyzed and the
  estimator $\frac{NL}{n}$ used, with value zero when $C = 0$. This preserves
  unbiasedness, but the variance formula~\ref{eq:nhat.variance} does not cover
  that fallback.
}

\begin{equation}
  H \mid C \sim \operatorname{Binomial}(m, q)
\end{equation}

which does not depend on $C$. Since $n$ and $m$ are just constants, this means
that $X$ and $Y$ are also independent, so~\ref{eq:rel.var} and \ref{eq:rel.ab}
are applicable. We have

\begin{align}
  a &= \frac{np(1-p)}{(np)^2} = \frac{1-p}{np}  \\
  b &= \frac{mq(1-q)}{(mq)^2} = \frac{1-q}{mq}
\end{align}

Using~\ref{eq:nhat}, \ref{eq:xy}, \ref{eq:rel.var}, \ref{eq:unbiased}

\begin{align}
  \Var(\hat{N}_{\text{legal}}) &= \Var(NXY) = N^2\Var(XY)  \\
    &= N^2 \E{XY}^2 (a + b + ab) = \E{NXY}^2 (a + b + ab)  \\
    &= \E{\hat{N}_{\text{legal}}}^2 (a + b + ab)           \\
    &= N_{\text{legal}}^2 (a + b + ab)
\end{align}

Finally, taking square roots:
\begin{equation}
  \label{eq:nhat.variance}
  \SD(\hat{N}_{\text{legal}}) = N_{\text{legal}} \sqrt{a + b + ab}
\end{equation}


% --------------------------------------------------------------------------------
\begin{thebibliography}{9}

\bibitem{wiki:binomial}
  \emph{Binomial distribution, Wikipedia, The Free Encyclopedia}, 2026.
  \url{https://en.wikipedia.org/wiki/Binomial_distribution}
      [Online; accessed 24-Sept-2026].

\bibitem{wiki:hypergeom}
  \emph{Hypergeometric distribution, Wikipedia, The Free Encyclopedia}, 2026.
  \url{https://en.wikipedia.org/wiki/Hypergeometric_distribution}
      [Online; accessed 24-Sept-2026].

\bibitem{wiki:totalexpect}
  \emph{Law of total expectation, Wikipedia, The Free Encyclopedia}, 2026.
  \url{https://en.wikipedia.org/wiki/Law_of_total_expectation}
      [Online; accessed 24-Sept-2026].

\bibitem{wiki:productvariance}
  \emph{Variance algebra for random variables, Wikipedia, The Free Encyclopedia}, 2026.
  \url{https://en.wikipedia.org/wiki/Algebra_of_random_variables#Variance_algebra_for_random_variables}
      [Online; accessed 24-Sept-2026].

\bibitem{wiki:bernoulli}
  \emph{Bernoulli distribution, Wikipedia, The Free Encyclopedia}, 2026.
  \url{https://en.wikipedia.org/wiki/Bernoulli_distribution}
      [Online; accessed 24-Sept-2026].

\end{thebibliography}

\end{document}
